% GENERATED FILE. Edit canonical lessons, then run npm run generate:books.
% canonical-source-hash: fnv1a64:10817f00f5743203
% canonical-lessons: MW-C182-jhatakno, MW-C182-gholno, MW-C182-milavno, MW-C182-bagharno, MW-C182-purasno

\chapter{Shake, Stir and Serve}
\label{ch:mw-shake-stir-and-serve}
\hlchaptermodality{182}

\begin{chapteropening}
\textbf{By the end of this chapter:} \emph{I can say how I shake out a sheet, stir, mix and serve.}

Complete the last lesson of chapter 182: I can say how I shake out a sheet, stir, mix and serve.
\end{chapteropening}

\section[\texorpdfstring{\mw{झटकणो} (jhaṭakṇo)}{झटकणो (jhaṭakṇo)}]{\mw{झटकणो} (jhaṭakṇo) — to shake out}

\label{lesson:MW-C182-jhatakno}



\begin{quote}
Before the new one: say the Marwadi for to strain, then the Marwadi for rice pudding.
\end{quote}

\subsection*{You'll want to know: \mw{झटकणो}}
\textbf{\mw{झटकणो}} — \emph{jhaṭakṇo} — \textquotedblleft{}to shake out\textquotedblright{}.

\begin{quote}
\mw{हूं} \mw{दिनूगै} \mw{चादर} \mw{झटकूं} \mw{हूं।}
\end{quote}

\begin{quote}
\emph{hū̃ dinūgai chādar jhaṭkū̃ hū̃.}
\end{quote}

\begin{quote}
I shake out the sheet early in the morning.
\end{quote}

Past (I …): \mw{म्हैं} \mw{चादर} \mw{झटकी}, \emph{mhaĩ chādar jhaṭkī}.

\begin{cousinweb}
\textbf{\mw{झटकणो}} (\emph{jhaṭakṇo}) comes from a reconstructed \emph{*jhaṭ}, \textquotedblleft{}a sudden movement\textquotedblright{}, the word behind \textbf{\mw{झट}}; CDIAL lists Hindi \emph{jhaṭaknā}, \textquotedblleft{}to jerk\textquotedblright{} (CDIAL 5327). CDIAL gives no Marwari form in this entry.
\end{cousinweb}

\begin{grammarlens}[title={What you've built}]
One more word for this chapter.
\end{grammarlens}

\subsection*{Your turn}
\begin{itemize}
\raggedright
  \item \emph{Say it:} \emph{jhaṭakṇo}
  \item \emph{Say it:} \emph{jhaṭakṇo}, once more
  \item \emph{Say it:} \emph{hū̃ dinūgai chādar jhaṭkū̃ hū̃.} — I shake out the sheet early in the morning.
  \item \emph{Recall:} say the Marwadi for to strain, then the Marwadi for rice pudding, then say \emph{jhaṭakṇo} again
\end{itemize}

\begin{tcolorbox}[breakable,colback=teal!4,colframe=teal!35!black,title={Before you move on}]
What does \mw{झटकणो} mean? (\textquotedblleft{}To shake out\textquotedblright{}.) Say it once more.
\end{tcolorbox}
\section[\texorpdfstring{\mw{घोळणो} (ghoḷṇo)}{घोळणो (ghoḷṇo)}]{\mw{घोळणो} (ghoḷṇo) — to dissolve; to stir in}

\label{lesson:MW-C182-gholno}



\begin{quote}
Before the new one: say the Marwadi for rice pudding, then the Marwadi for to shake out.
\end{quote}

\subsection*{You'll want to know: \mw{घोळणो}}
\textbf{\mw{घोळणो}} — \emph{ghoḷṇo} — \textquotedblleft{}to dissolve; to stir in\textquotedblright{}.

\begin{quote}
\mw{हूं} \mw{पाणी} \mw{मांय} \mw{गुड़} \mw{घोळूं} \mw{हूं।}
\end{quote}

\begin{quote}
\emph{hū̃ pāṇī mā̃y guṛ ghoḷū̃ hū̃.}
\end{quote}

\begin{quote}
I dissolve jaggery in water.
\end{quote}

Past (I …): \mw{म्हैं} \mw{गुड़} \mw{घोळ्यो}, \emph{mhaĩ guṛ ghoḷyo}.

\begin{cousinweb}
\textbf{\mw{घोळणो}} (\emph{ghoḷṇo}) is inherited from Sanskrit \emph{gholayati}, \textquotedblleft{}mixes, stirs together\textquotedblright{}; CDIAL lists Hindi \emph{gholnā}, \textquotedblleft{}to mix, dissolve\textquotedblright{}, and Gujarati \emph{ghoḷvũ} (CDIAL 4526). CDIAL gives no Marwari form in this entry.
\end{cousinweb}

\begin{grammarlens}[title={What you've built}]
One more word for this chapter.
\end{grammarlens}

\subsection*{Your turn}
\begin{itemize}
\raggedright
  \item \emph{Say it:} \emph{ghoḷṇo}
  \item \emph{Say it:} \emph{ghoḷṇo}, once more
  \item \emph{Say it:} \emph{hū̃ pāṇī mā̃y guṛ ghoḷū̃ hū̃.} — I dissolve jaggery in water.
  \item \emph{Recall:} say the Marwadi for rice pudding, then the Marwadi for to shake out, then say \emph{ghoḷṇo} again
\end{itemize}

\begin{tcolorbox}[breakable,colback=teal!4,colframe=teal!35!black,title={Before you move on}]
What does \mw{घोळणो} mean? (\textquotedblleft{}To dissolve; to stir in\textquotedblright{}.) Say it once more.
\end{tcolorbox}
\section[\texorpdfstring{\mw{मिलावणो} (milāvṇo)}{मिलावणो (milāvṇo)}]{\mw{मिलावणो} (milāvṇo) — to mix}

\label{lesson:MW-C182-milavno}



\begin{quote}
Before the new one: say the Marwadi for to shake out, then the Marwadi for to dissolve.
\end{quote}

\subsection*{You'll want to know: \mw{मिलावणो}}
\textbf{\mw{मिलावणो}} — \emph{milāvṇo} — \textquotedblleft{}to mix\textquotedblright{}.

\begin{quote}
\mw{हूं} \mw{दूध} \mw{मांय} \mw{खांड} \mw{मिलाऊं} \mw{हूं।}
\end{quote}

\begin{quote}
\emph{hū̃ dūdh mā̃y khā̃ḍ milāū̃ hū̃.}
\end{quote}

\begin{quote}
I mix sugar into the milk.
\end{quote}

Past (I …): \mw{म्हैं} \mw{दूध} \mw{मांय} \mw{खांड} \mw{मिलाई}, \emph{mhaĩ dūdh mā̃y khā̃ḍ milāī}.

\begin{cousinweb}
\textbf{\mw{मिलावणो}} (\emph{milāvṇo}) is built on Sanskrit \emph{milati}, \textquotedblleft{}meets\textquotedblright{}; CDIAL lists Old Marwari \emph{milaï} (CDIAL 10133). It is the causative partner of \textbf{\mw{मिलणो}}, to meet: to make things meet is to mix them.
\end{cousinweb}

\begin{grammarlens}[title={What you've built}]
One more word for this chapter.
\end{grammarlens}

\subsection*{Your turn}
\begin{itemize}
\raggedright
  \item \emph{Say it:} \emph{milāvṇo}
  \item \emph{Say it:} \emph{milāvṇo}, once more
  \item \emph{Say it:} \emph{hū̃ dūdh mā̃y khā̃ḍ milāū̃ hū̃.} — I mix sugar into the milk.
  \item \emph{Recall:} say the Marwadi for to shake out, then the Marwadi for to dissolve, then say \emph{milāvṇo} again
\end{itemize}

\begin{tcolorbox}[breakable,colback=teal!4,colframe=teal!35!black,title={Before you move on}]
What does \mw{मिलावणो} mean? (\textquotedblleft{}To mix\textquotedblright{}.) Say it once more.
\end{tcolorbox}
\section[\texorpdfstring{\mw{बघारणो} (baghārṇo)}{बघारणो (baghārṇo)}]{\mw{बघारणो} (baghārṇo) — to temper (spices in hot oil)}

\label{lesson:MW-C182-bagharno}



\begin{quote}
Before the new one: say the Marwadi for to dissolve, then the Marwadi for to mix.
\end{quote}

\subsection*{You'll want to know: \mw{बघारणो}}
\textbf{\mw{बघारणो}} — \emph{baghārṇo} — \textquotedblleft{}to temper (spices in hot oil)\textquotedblright{}.

\begin{quote}
\mw{मां} \mw{दाल} \mw{बघारै} \mw{है।}
\end{quote}

\begin{quote}
\emph{mā̃ dāl baghārai hai.}
\end{quote}

\begin{quote}
Mother tempers the dal.
\end{quote}

Past (I …): \mw{म्हैं} \mw{दाल} \mw{बघारी}, \emph{mhaĩ dāl baghārī}.

\begin{cousinweb}
\textbf{\mw{बघारणो}} (\emph{baghārṇo}) is inherited from Sanskrit \emph{vyāghārayati}, \textquotedblleft{}sprinkles over\textquotedblright{} (with ghee); CDIAL lists Hindi \emph{baghārnā}, for frying spices in hot oil (CDIAL 12191). CDIAL gives no Marwari form in this entry.
\end{cousinweb}

\begin{grammarlens}[title={What you've built}]
One more word for this chapter.
\end{grammarlens}

\subsection*{Your turn}
\begin{itemize}
\raggedright
  \item \emph{Say it:} \emph{baghārṇo}
  \item \emph{Say it:} \emph{baghārṇo}, once more
  \item \emph{Say it:} \emph{mā̃ dāl baghārai hai.} — Mother tempers the dal.
  \item \emph{Recall:} say the Marwadi for to dissolve, then the Marwadi for to mix, then say \emph{baghārṇo} again
\end{itemize}

\begin{tcolorbox}[breakable,colback=teal!4,colframe=teal!35!black,title={Before you move on}]
What does \mw{बघारणो} mean? (\textquotedblleft{}To temper (spices in hot oil)\textquotedblright{}.) Say it once more.
\end{tcolorbox}
\section[\texorpdfstring{\mw{पुरसणो} (purasṇo)}{पुरसणो (purasṇo)}]{\mw{पुरसणो} (purasṇo) — to serve food}

\label{lesson:MW-C182-purasno}



\begin{quote}
Before the new one: say the Marwadi for to mix, then the Marwadi for to temper.
\end{quote}

\subsection*{You'll want to know: \mw{पुरसणो}}
\textbf{\mw{पुरसणो}} — \emph{purasṇo} — \textquotedblleft{}to serve food\textquotedblright{}.

\begin{quote}
\mw{मां} \mw{थाळी} \mw{मांय} \mw{रोटी} \mw{पुरसै} \mw{है।}
\end{quote}

\begin{quote}
\emph{mā̃ thāḷī mā̃y roṭī pursai hai.}
\end{quote}

\begin{quote}
Mother serves rotis onto the plate.
\end{quote}

Past (I …): \mw{म्हैं} \mw{खीचड़ो} \mw{पुरस्यो}, \emph{mhaĩ khīchṛo purasyo}.

\begin{cousinweb}
\textbf{\mw{पुरसणो}} (\emph{purasṇo}) is inherited from Sanskrit \emph{pariveṣati}, \textquotedblleft{}serves food\textquotedblright{}; CDIAL lists Old Marwari \emph{purasaï}, \textquotedblleft{}serves a meal\textquotedblright{} (CDIAL 7888).
\end{cousinweb}

\begin{grammarlens}[title={What you've built}]
One more word for this chapter.
\end{grammarlens}

\subsection*{Your turn}
\begin{itemize}
\raggedright
  \item \emph{Say it:} \emph{purasṇo}
  \item \emph{Say it:} \emph{purasṇo}, once more
  \item \emph{Say it:} \emph{mā̃ thāḷī mā̃y roṭī pursai hai.} — Mother serves rotis onto the plate.
  \item \emph{Recall:} say the Marwadi for to mix, then the Marwadi for to temper, then say \emph{purasṇo} again
\end{itemize}

\begin{tcolorbox}[breakable,colback=teal!4,colframe=teal!35!black,title={Before you move on}]
What does \mw{पुरसणो} mean? (\textquotedblleft{}To serve food\textquotedblright{}.) Say it once more.
\end{tcolorbox}
